\section[q6]{Question 6}
\textit{We are concerned that population size might modify the relationship between who\_region
and the highest air pollution exposure category you created. Determine if there is
any evidence of effect modification.}

\subsection{Question 6.a}
\textit{Create either a single table or multiple tables to summarise the relationship
between who\_region and highest air pollution exposure category you cre-
ated, grouped by population size category you created in Assignment 1.
If you have more than 3 population size categories, use the first three categories
for this analysis.}

\subsection{Question 6.b.i}
\textit{Investigate whether the relationship between who\_region and highest air pol-
lution exposure category you created varies by population size category you
created in Assignment 1. Be sure to specify which groups (4\_Eur or 1\_Afr)
you are treating as the exposed and unexposed groups. Your outcome should
be highest air pollution exposure category you created.}

\subsubsection{Question 6.b.ii}
\textit{To do this, calculate the prevalence ratio in the three population size cat-
egories (you can do this by hand or in your software).
Report whether the association between who\_region and highest air
pollution exposure category you created is statistically significant in each
of the three population size categories by using 95\% confidence intervals
or by conducting chi-squared tests of the association in each of the three
groups. Write out your answer for each group in a sentence (do not just
report the numbers).}


\subsubsection{Question 6.b}
\textit{Does population size appear to be modifying the effect of who\_region
on highest air pollution exposure category you created? Justify your an-
swer using the information and results calculated above. (3)
Biostatistics 1: Assignment 2 5 of 6}

\subsection{Question 6.c}
\textit{Based on your analysis and interpretation above, select which of the results
from the list below you would report in your manuscript and why.}

\subsubsection{Question 6.c.i}
\textit{Just the crude overall prevalence ratio.}

\subsubsection{Question 6.c.ii}
\textit{The stratified prevalence ratios (in each population category).}

\subsubsection{Question 6.c.iii}
\textit{An adjusted prevalence ratio like the Mantel Haenszel adjusted association (adjusted for population category).
Hint: You need to filter your data to include just one population category at a time then do the next population category to do these grouped
analyses
}


\newpage